\documentclass[11pt,a4paper]{article}
\usepackage[margin=23mm]{geometry}
\usepackage{amsmath,amssymb,booktabs}
\usepackage[hidelinks]{hyperref}
\usepackage[T1]{fontenc}
\setlength{\parindent}{0pt}
\setlength{\parskip}{6pt}
\begin{document}
\begin{center}
{\Large VSL-B1: information and departure-time examples}\\[4pt]
13 October 2026 --- original conditional deductions, not an adopted model
\end{center}

\section*{1. Combined continuing mass need not determine adjacent receiving}
Consider equal-length local sections in three lanes: ending $E$, adjacent
continuing $A$, and farther continuing $F$. Length is $\ell$ [m]; densities
$k_E,k_A,k_F$ are in vehicle-equivalents/m. For this diagnostic only, use the
standard receiving relation
\[
R(k)=\min\{C,\;w(k_j-k)\},
\]
where $C$ is a rate [vehicle-equivalents/s], $w>0$ is a backward-wave speed
[m/s], and $k_j$ is jam density. This is a source-compatible receiving
ingredient, not a selected VSL or lane-transfer update.

Choose $0\leq k_L<k_H\leq k_j$ with $R(k_H)<R(k_L)$. The two snapshots
\[
 (k_E,k_A,k_F)=(k_E,k_H,k_L),\qquad
 (k_E,k_A,k_F)=(k_E,k_L,k_H)
\]
have the same ending mass $\ell k_E$ and combined continuing mass
$\ell(k_H+k_L)$, but different receiving opportunity in the adjacent lane.
With identical positive merge requests and an allocation in which that
receiving difference binds, completed transfer can differ.

\textbf{Conclusion.} No function of ending and combined continuing mass alone
can recover adjacent receiving for every snapshot in this set. This does not
establish a reachable equal-full-history pair, a different optimal VSL, or
necessity of a full lane profile. A restricted domain or compact receiving
statistic may remove the ambiguity.

\textbf{Initial symmetry is a separate issue.} Let $a,f$ be adjacent and farther
continuing masses and let $h>0$ be a transfer rate from $E$ to $A$. If the two
continuing lanes have equal net longitudinal contributions $b$ at a comparison
instant and no compensating redistribution, then
\[
 \dot a=b+h,\qquad \dot f=b,\qquad
 \frac{d}{dt}(a-f)=h>0.
\]
Thus $a=f$ initially does not imply continued equality. A dynamically valid
grouped model needs an explicit equivalence condition. These equations are
local conservation diagnostics, not proposed traffic evolution laws.

\vfill
Primary receiving lineage: Carey--Balijepalli--Watling, Sections 3--5.\\
\href{https://eprints.whiterose.ac.uk/id/eprint/77315/7/Carey\%20et\%20al\%202013_with_coversheet.pdf}{institutional manuscript};
FB0030, equations 13--15, archived in the review package.

\newpage
\section*{2. Equal final work can give different mean completion times}
This example examines the exit interface only. All class demand is exposed at
time zero. Initial continuing mass is $c_0=1$ and merge mass is $e_0=1$
[vehicle-equivalents]. Ordinary service costs one resource unit per unit mass;
merge service costs $\alpha=2$. Constant receiving work rate is
$\mu=1$ [resource units/s]. Mass is continuous; no integer vehicle process is
being introduced.

Every work-conserving allocation uses total work
\[
 W=c_0+\alpha e_0=3,
 \qquad T=W/\mu=3\ \text{s}.
\]
Let $N(t)$ be total mass not yet exited. Complete population time is
$F=\int_0^T N(t)\,dt$ [vehicle-equivalent seconds]; mean time is
$\bar T=F/(c_0+e_0)$.

\textbf{Continuing first.} Ordinary traffic clears during the first second;
merge traffic then clears over two seconds:
\[
 N_{C}(t)=
 \begin{cases}2-t,&0\leq t\leq1,\\(3-t)/2,&1\leq t\leq3.\end{cases}
 \qquad F_C=\frac52,\quad \bar T_C=\frac54\ \text{s}.
\]
\textbf{Merge first.} Merge traffic clears during the first two seconds;
ordinary traffic then clears over one second:
\[
 N_{E}(t)=
 \begin{cases}2-t/2,&0\leq t\leq2,\\3-t,&2\leq t\leq3.\end{cases}
 \qquad F_E=\frac72,\quad \bar T_E=\frac74\ \text{s}.
\]
\textbf{Proportional co-clearance.} S13's conditional proportional resource
projection preserves the class ratio, giving
\[
 N_P(t)=2\left(1-\frac{t}{3}\right),\quad
 F_P=3,\quad \bar T_P=\frac32\ \text{s}=T/2.
\]
\begin{center}
\begin{tabular}{lccc}
\toprule
Allocation & Final work & Final clearance [s] & Mean time [s]\\
\midrule
Continuing first & 3 & 3 & 1.25\\
Merge first & 3 & 3 & 1.75\\
Proportional & 3 & 3 & 1.50\\
\bottomrule
\end{tabular}
\end{center}
\textbf{Conclusion.} Final work or clearance time alone does not determine the
complete-time objective under general priorities. This is a hand-derived
counterexample, not a result attributed to a source experiment. Source priority
alternatives motivate checking allocation; they do not select these schedules
as the proposed freeway model.

\newpage
\section*{3. What these examples do and do not qualify}
The examples establish two interface risks: loss of consequential receiver
information and loss of consequential departure timing. Neither selects a
replacement lane-transfer model, validates its common-speed response, or proves
an information-to-optimal-action split.

For the fixed finite cohort with mass $M$, count the mass $N_{\mathrm{sys}}(t)$
still at sources or in the modeled system. With a common decision/accounting
origin, complete exit at $T(v)$ gives
\[
 F(v)=\int_0^{T(v)}N_{\mathrm{sys}}(t;v)\,dt,
 \qquad \bar T(v)=F(v)/M.
\]
If admitted trips have later release times, subtract their fixed release-time
prefix in interpreting trip elapsed times. It is action-independent; it may
not be confused with omitted action-dependent source waiting. No new demand
release regime is selected here. The objective meaning remains the accepted
complete population-time meaning even if its reduced formula changes.

\textbf{Ideal travel is conditional.} $L/v$ is a traversal time only when the
declared common advice $v$ actually determines the relevant progress over
length $L$. Binding receiving can invalidate this ideal mapping. The review
does not supply an alternative speed/fundamental-diagram law.

\textbf{MP6 is a joint qualification.} One admitted regime must support the
same geometry, finite cohort, advice response, receiving and grouping,
early/late cost accounting, lane-end boundary and departures throughout its
declared action range. Source-compatible ingredients and correct algebra alone
do not establish that regime, calibration, or useful external gradients.

\textbf{Minimal-state and comparison limits.} S13's autonomous initial-ending-
mass scalar remains exactly sufficient under S13. The receiver diagnostic
does not invalidate that result; it concerns a different state domain. A new
candidate must compare against a compact predictor given the same admissible
initial information. Native reachable history matching and independent
perturbation usefulness are still unproved.

\textbf{Current gate.} VSL-B1-G1 asks whether qualification must preserve the
restricted sequential scalar/batch reduction (B1-R), or may propose the
smallest local receiving and within-cohort timing amendment (B1-L). No source
law or new state is adopted by this note. No simulation or parameter sweep
was run.

\vfill
Companion files in the review package:\\
\texttt{vsl\_b1\_lane\_change\_model\_and\_gate.md}\\
\texttt{vsl\_b1\_source\_and\_interface\_audit.md}.
\end{document}
